% BANKS-SOURCE: pdf=155 print=145 kind=section id=9.3
\subsection{Renormalization and critical phenomena}

% BANKS-SOURCE: pdf=155 print=145 kind=prose
The deep understanding of renormalization and the renormalization group came
from the work of Fisher, Kadanoff, Wilson, and other condensed-matter
physicists, working on the problem of continuous phase transitions, or
\emph{critical phenomena}. We all know that many forms of ordinary matter
undergo phase transitions as we vary control parameters like pressure,
temperature, and external magnetic fields. Water turns into steam, solids
melt, ferromagnets magnetize, etc. etc. The thermodynamics of a system is
characterized by a free energy per unit volume,
$F(T,P,H,\ldots)$, which depends on the control parameters. Typically $F$
depends analytically on these parameters, and, for systems with a finite number
of degrees of freedom, we can prove that this is so. Mathematically, phase
transitions occur only for infinite volume (unless we have an infinite number
of fields in our system at each point), and show up as discontinuities in
derivatives of $F$ w.r.t. e.g. the temperature. Ehrenfest defined the order of
a transition as the order of the derivative which became singular or
discontinuous at the phase transition.

\BanksImplicitHook{I-CH09-004}

Landau, who invented the first general theory of phase transitions
(generalizing van der Waals' treatment of the liquid--gas system), realized
that the main difference was between first-order transitions and all the
others, which are called continuous transitions. Landau's theory is extremely
simple. He postulated that the different phases of the system were
characterized by an \emph{order parameter} $\Phi$ distinguishing the phases.
For example, in a ferromagnet $\Phi$ is the magnetization. In the liquid--gas
transition, it is related to the difference between the density of the phases
and the critical density at which the transition occurs.

The order parameter is defined so that its value is small near the transition.
It is a macroscopic property of the system, so it does not have thermal
fluctuations. Its value is therefore determined by minimizing a free-energy
function $F(\Phi,T,\ldots)$ w.r.t. $\Phi$, at fixed values of the control
parameters. Near the transition, we can expand
\[
  F=a(T)\Phi^2+b(T)\Phi^3+c(T)\Phi^4+\cdots .
\]
The coefficients in the expansion depend analytically on the control
parameters, of which we have indicated only the temperature. Keeping only the
indicated terms, and assuming that $c>0$, we have a quartic polynomial, with
two local minima. As we vary $T$, it might happen that the global minimum
jumps from one to the other. This is a first-order phase transition: $F$ is
continuous, but its first derivative is not, because we have jumped from one
branch of the function $\Phi(T)$ to another when the free energies of the two
branches cross each other.

Landau then realized that one could explain second-order transitions by
assuming that $b$ was zero (either by fine tuning some other control
parameter, or because the system had a built in $\Phi\to-\Phi$ symmetry) and
that, at the critical temperature, $a(T_c)=0$. At this point two minima
coalesce into one. The magnetization and free energies have square-root branch
points in $T-T_c$.

% BANKS-SOURCE: pdf=156 print=146 kind=prose
Landau's theory, called \emph{mean field theory}, has two remarkable
properties: it is completely universal, depending only on the properties of
polynomials, not on the detailed nature of the order parameter or the system it
characterizes; and it predicts fractional power-law exponents for physical
properties near continuous phase transitions. Experimentalists rushed to test
these properties. The experiments were difficult, because one had to get very
near the critical point to measure the exponents.

The result of this experimental work was exciting and puzzling. Continuous
phase transitions did exhibit remarkable universality. For example, the
critical point in the liquid--vapor phase diagram in the $P,T$ plane has the
same critical exponents as the so-called Ising ferromagnets (ferromagnets with
a single preferred axis of up--down polarization). However, there was not
complete universality. For example, XY ferromagnets with a whole preferred
plane of polarization, or those in which there is no preferred axis, had
different critical exponents. Furthermore, while the exponents were not
wildly different from the Landau square roots, they were definitely different.
Clearly something was right about Landau's ideas, but something much more
intricate and interesting was going on.

This was confirmed by another set of phenomena,\footnote[5]{The most
spectacular of which is \emph{critical opalescence}.} which were observed by
scattering light and neutrons off of substances undergoing second-order
transitions. These phenomena could be explained by assuming that correlation
functions of local quantities, like the local magnetization, obeyed scaling
laws of the form
\[
  \langle M(\mathbf x)M(\mathbf y)\rangle\sim|\mathbf x-\mathbf y|^{-2\beta},
\]
at very large distances and temperatures at the critical point. By contrast,
away from critical points, three-dimensional condensed-matter systems
generally exhibit what is called a \emph{correlation length}
\[
  \langle M(\mathbf x)M(\mathbf y)\rangle\sim
  \frac{\mathrm{e}^{-|\mathbf x-\mathbf y|/l_c}}{|\mathbf x-\mathbf y|} .
\]
The three new phenomena of \emph{universality classes}, rather than a
completely universal \emph{mean field theory}, \emph{non-mean field critical
exponents}, and \emph{infinite correlation length} at the critical point,
received a common explanation in terms of the renormalization group.

Kadanoff's first crude form of the renormalization group is called the
\emph{block spin transformation}. Kadanoff first reasoned that, if critical
phenomena had universal aspects, then interesting results for real systems
could be extracted from simple models like the Ising model. This model
consists of a classical spin, which can take only the values $\pm1$, situated
on the sites of a $d$-dimensional hypercubic lattice (where $d=1,2$ or 3 for
systems realizable in condensed-matter laboratories). The Hamiltonian is
\[
  H_0=-J\sum_{\mathbf x,\mu}\sigma(\mathbf x)\sigma(\mathbf x+\mathbf e_\mu),
\]
where $\mathbf e_\mu$ is a lattice unit vector and the sum runs over all
nearest-neighbor points in the lattice.
Kadanoff's great insight was to realize that, if critical phenomena really had
something to do with an infinite correlation length, then the microscopic
dynamics was

% BANKS-SOURCE: pdf=157 print=147 kind=prose
irrelevant and there should be some sort of long-wavelength effective theory
that captured the phenomena. Landau's theory could be viewed as a sort of
guess at what this long-wavelength theory was (a so-called Gaussian model), but
the experiments showed that this guess was not correct. How should one
systematically calculate the effective long-wavelength Hamiltonian?

The slogan that works, as it does for so many other things, is \emph{little
steps for little feet}. If you're working on a lattice with spacing $a$, and
you want to investigate phenomena on a length scale much bigger than $a$,
first go to $2a$, then from there to $4a$, and so on. Kadanoff did this by
defining blocks on the lattice, and defining $\Phi_i$ to be the average of the
spins on the $i$th block. He imagined doing the partition sum over all spins,
with the $\Phi_i$ held fixed. One then gets a new system, whose variables are
the $\Phi_i$, defined on a lattice of spacing $2a$. The Hamiltonian takes the
form
\[
  H_1=\sum K_n^1(\mathbf x_1,\ldots,\mathbf x_n)
  \Phi(\mathbf x_1)\ldots\Phi(\mathbf x_n),
\]
where the sums now go over all lattice points and all $n$.

It isn't easy to calculate the functions $K_n^1$, but it is easy to argue that
they all have short-range correlations, because the original Hamiltonian did,
and, by keeping the $\Phi_i$ fixed, we have not allowed any long-wavelength
fluctuations. That is, the sum over $\Phi_i$, which we have not yet done,
contains both configurations where the $\Phi_i$ randomly fluctuate from point
to point on the lattice with spacing $2a$ and $\Phi_i$ configurations of
longer wavelength. It is only by summing over the latter that we can produce
long-range correlations in the system. We can now iterate the procedure to get
new effective theories (with the same partition function) for lattices of
larger and larger spacing. After $k$ steps we have
\[
  H_k=\sum K_n^k(\mathbf x_1,\ldots,\mathbf x_n)
  \Phi(\mathbf x_1)\ldots\Phi(\mathbf x_n),
\]
where the functions $K_n^k$ have short-range correlations on a lattice of
spacing $2^ka$. Note that we have used the same name for the variables at all
steps but the first. Actually these variables all have different discrete
ranges, but as $k$ gets large it is clear that we can let the variables be
continuous and replace sums over $\Phi$ by integrals. Furthermore, although
all the $\Phi$ variables are bounded between $\pm1$ we can abandon this
constraint and replace it by a potential that favors $\Phi$s in this compact
range. Our system now resembles a lattice field theory, with the
statistical-mechanical Hamiltonian playing the role of the Euclidean action.

\BanksImplicitHook{I-CH09-005}

As we take $k$ to infinity, one of two things happens. If we are at a value of
the temperature for which the correlation length is finite, then the functions
$K_n^k$ become of shorter and shorter range in lattice units, once we take the
spacing larger than the correlation length. Eventually, they are all
Kronecker deltas, and there are no interactions between the lattice points. We
call this limit the \emph{trivial fixed point} of the renormalization group.
But suppose that we are at the critical temperature, for which the correlation
length is infinite. In this case, Kadanoff realized that the only way to get a
sensible limit was to assume that the Hamiltonian approached a fixed point, in
terms of appropriately rescaled variables $\hat\Phi$. That is, the fixed-point
Hamiltonian had to be invariant under

% BANKS-SOURCE: pdf=158 print=148 kind=page-boundary
scale transformations of space combined with a rescaling of $\Phi$. This then
predicts scale-invariant correlation functions, as observed experimentally. It
further predicts that the scaling exponents are characteristic of the
fixed-point Hamiltonian and independent of the microphysics. This explains
universality.

However, unlike Landau's theory, Kadanoff's block spin idea can accommodate
different universality classes. For example, there is no reason for the RG
transformation for a system with Ising symmetry to be the same as that for
planar or fully rotation-invariant spin symmetry. Detailed calculation shows
that it is not. Thus, block spin or renormalization theory explains the
qualitative nature of critical phenomena. It remained for Fisher, Wilson, and
others to appreciate the connection to Euclidean field theory, and to use
field-theoretic tools to calculate detailed numerical agreement of the theory
with a wide variety of experiments. Rather than following them on this
fascinating journey, we will now return to the application of these ideas to
relativistic quantum field theory itself.
